\documentclass[10pt,a4paper]{amsart}
\usepackage[margin=19mm]{geometry}
\usepackage{amsmath,amssymb,mathtools}
\usepackage[scr=rsfs,cal=euler]{mathalfa}
\usepackage{xfrac}
\usepackage[unicode,hidelinks,bookmarks=false]{hyperref}
\newcommand{\ie}{i.e.\ }
\newcommand{\cf}{cf.\ }
\newcommand{\resp}{respectively\ }
\newcommand{\Subsection}{\subsection}
\providecommand{\nobreakdash}{\nobreak\hbox{-}}
\setlength{\emergencystretch}{2em}
\multlinegap0pt
\usepackage{amssymb}
\usepackage{xcolor}
\usepackage{float}
\usepackage{tikz}
\newtheorem{prop}{Proposition}
\newtheorem{theo}{Theorem}
\newtheorem{coro}{Corollary}
\newtheorem{lemm}{Lemma}
\def\Gal{\mathrm{Gal}}
\def\Pic{\mathrm{Pic}}
\def\bP{{\mathbb P}}
\def\bQ{{\mathbb Q}}
\def\bZ{{\mathbb Z}}
\def\cT{{\mathcal T}}
\def\disc{{\mathrm{disc}}}
\def\fC{{\mathfrak C}}
\def\fD{{\mathfrak D}}
\def\fS{{\mathfrak S}}
\def\rH{{\mathrm H}}
\def\sD{{\mathsf D}}
\title{Universal torsors over quartic del Pezzo surfaces and stable rationality}
\author{Yuri Tschinkel, Zhijia Zhang}
\date{Source: arXiv:2608.20029v2 (CC BY 4.0)}
\begin{document}
\maketitle

\noindent\emph{Source and licence.} Universal torsors over quartic del Pezzo surfaces and stable rationality. Version arXiv:2608.20029v2 (CC BY 4.0), distributed by arXiv under the Creative Commons Attribution 4.0 International licence, http://creativecommons.org/licenses/by/4.0/, as recorded in the arXiv licence metadata for this exact version and shown on the abstract page https://arxiv.org/abs/2608.20029v2. Full licence notice: \texttt{licenses/arxiv-2608.20029-CC-BY-4.0.txt}.\par\smallskip

\noindent\textit{Editorial scope.} Counted unit: the article's prop prop:cubictypeI3 (source0.tex lines 1050--1057). The article's complete original proof of it is delivered with it (source0.tex lines 1058--1177, one \texttt{proof} environment of 120 lines). No mathematics is written, repaired or re-derived: the statement and the proof are the article's own lines, in the article's own notation, and no retained line is substituted or re-flowed.  Retained uncounted as the source-local context the proof needs: lines 295--298; lines 470--478; lines 683--685; lines 706--711; lines 764--776; lines 779--879.  Nothing was omitted from any retained span, so the package carries no omission record.  The retained lines cite 6 works, whose entries are transcribed verbatim from the article's own bibliography source in the e-print and delivered with short editorial labels recorded in the item's reference mapping.  No retained line contains a figure environment, an image-inclusion command or drawing code, so no image is delivered and the package declares no asset dependency.  Editorial formatting only, in addition: an AMS \texttt{amsart} preamble that restates the article's own declarations and macros used by the retained lines, namely the environments \texttt{prop}, \texttt{theo}, \texttt{coro}, \texttt{lemm} on separate counters, together with the standard packages those lines need.  The retained environments print in this document as Proposition 1, Theo 1, Coro 1, Lemm 1 in order of appearance, whereas the article prints the same items with the numbering of its own full document.  The complete native source of the article as distributed in the arXiv e-print for this exact version is bundled unchanged as \texttt{sources/208148/source0.tex}.
    \begin{equation} 
    \label{eqn:h1}
\rH^1(\mathfrak g_{k'}, \Pic(\bar{S}))=0, \quad \text{ for all finite }  k'/k. 
    \end{equation}

\begin{prop}[{\cite[Proposition 3]{bctspstably}}]\label{prop:bctsmain}
    Let 
$X$ be a smooth, proper, geometrically integral, and geometrically rational 
variety over $k$. Assume that 
there exists a universal torsor 
$\cT\to X$ which is a
$k$-rational variety, and the 
$\mathfrak g$-module $\Pic(\bar{X})$ is a stably permutation module. Then $X$ is stably rational over $k$.
\end{prop}

\begin{theo}\label{thm:mainrepeat}
    Let $S$ be a quartic del Pezzo surface over $k$ and $\cT\to S$ a universal torsor over $S$. If $\cT(k)\neq \varnothing$, then  $\cT$ is rational over $k$. 
\end{theo}

\begin{coro}\label{coro:rattorsor}
    Let $S$ be a quartic del Pezzo surface over $k$ such that $S(k)\ne \varnothing$. Then there exists a $k$-rational universal torsor $\cT\to S$.
\end{coro}
\begin{proof}
    By Theorem~\ref{thm:mainrepeat}, it suffices to show the existence of a universal torsor $\cT\to S$ with a rational point. Recall that $S(k)\ne \varnothing$ implies that we can find $s\in S(k)$ in the complement to exceptional curves. By \cite[Corollary 2.3.4]{CTSansucDuke}, there exists a universal torsor $\cT$ over $S$ whose fiber over $s$ is a trivial torsor, and thus $\cT(k)\ne \varnothing$.
\end{proof}

\begin{prop}[{\cite{KSS}, \cite{TY}}]\label{prop:h1}
    If a quartic del Pezzo surface $S$ is $k$-minimal, and satisfies Condition {\bf (H1)} of \eqref{eqn:h1}, i.e., 
    $$
    \rH^1(\mathfrak g_{k'},\Pic(\bar{S}))=0,\quad \text{for all finite } k'/k,
    $$
   then the image of $\mathfrak g_{k}$ in $W(\sD_5)$ is conjugate to one of the following:
   \begin{itemize}
       \item[$(I_0):$] $\fS_3=\langle (3,4,5),c_1c_3c_4c_5(4,5)\rangle$,
         \item[$(I_1):$] $\fC_2\times\fS_3=\langle c_2c_3(4,5),c_1c_2c_3c_4(3,4,5)\rangle$,
           \item[$(I_2):$] $\fC_3\rtimes\fC_4=\langle (1,2,3),c_1c_2c_3c_4(2,3)(4,5)\rangle$,
             \item[$(I_3):$] $\fC_3\rtimes\fD_4=\langle c_1c_2c_3c_5(2,5),c_3c_4(3,4)(1,5,2)\rangle$.
   \end{itemize}
   \end{prop}

\begin{lemm}\label{lemm:stablyperm}
    Let $S$ be a quartic del Pezzo surface over $k$. Then 
    the $\mathfrak g$-module 
    $\Pic(\bar{S})$ is stably permutation if and only if Condition {\bf (H1)} of  \eqref{eqn:h1} is satisfied.
\end{lemm}
   \begin{proof}
      If $S$ is not $k$-minimal, the assertion follows from 
      %the fact that a del Pezzo surface of degree $\geq 5$ with a $k$-point is $k$-rational 
      \cite{manin-book}. Thus we may assume that $S$ is $k$-minimal. It suffices to show that $\Pic(\bar{S})$ is stably permutation, for the four types in Proposition~\ref{prop:h1}. The type $I_0$ is addressed in \cite{bctspstably}. Up to conjugation in $W(\sD_5)$, the groups of type $I_0$, $I_1$ and $I_2$ are subgroups of the group of type $I_3$. Thus we only need to consider this last type. 
      Let 
      $$
      g_1=c_1c_2c_3c_5(2,5), \quad g_2=c_3c_4(3,4)(1,5,2).
      $$
      In the standard basis $H, E_1,\ldots,E_5$ of  $\Pic(\bar{S})$, they act by matrices
{\small
$$
g_1=
\begin{pmatrix}
 3& 1& 1& 1& 2& 1\\
-1&-1& 0& 0&-1& 0\\
-1& 0& 0& 0&-1&-1\\
-1& 0& 0&-1&-1& 0\\
-2&-1&-1&-1&-1&-1\\
-1& 0&-1& 0&-1& 0
\end{pmatrix},
\quad
g_2=
\begin{pmatrix}
 2& 1& 1& 0& 0& 1\\
-1&-1& 0& 0& 0&-1\\
-1&-1&-1& 0& 0& 0\\
 0& 0& 0& 1& 0& 0\\
 0& 0& 0& 0& 1& 0\\
-1& 0&-1& 0& 0&-1
\end{pmatrix}.
$$}
In our convention, the matrices act from the right, e.g.,  
$$
g_2(H)=2H-E_1-E_2-E_5.
$$
We introduce new variables $w_0,w_1,w_2,q_0,q_1$ with $\mathfrak g$-permutation actions generated by 
$$
g_1(w_0)=w_0,\quad g_1(w_1)=w_2,\quad g_1(w_2)=w_1,\quad g_1(q_0)=q_1,\quad g_1(q_1)=q_0,
$$
$$
g_2(w_0)=w_2,\quad g_2(w_1)=w_0,\quad g_2(w_2)=w_1,\quad g_2(q_0)=q_0,\quad g_2(q_1)=q_1.
$$
One can check that 
$$
\Pic(\bar{S})\oplus\bZ[w_0,w_1,w_2,q_0,q_1]
$$ 
is a permutation module. Indeed, in the basis
\begin{align*}
b_1={}&
  2H-E_2-E_3-E_4-E_5-w_0-w_2-2q_0-q_1,\\
b_2={}&
  H-E_1-w_0-w_1-q_0-2q_1,\\
b_3={}&
  2H-E_1-E_3-E_4-E_5-w_0-w_1-2q_0-q_1,\\
b_4={}&
  H-E_2-w_1-w_2-q_0-2q_1,\\
b_5={}&
  H-E_5-w_0-w_2-q_0-2q_1,\\
b_6={}&
  2H-E_1-E_2-E_3-E_4-w_1-w_2-2q_0-q_1,\\
b_7={}&
  -3H+E_1+E_2+2E_3+E_4+E_5
  +w_0+w_1+w_2+3q_0+q_1,\\
b_8={}&
  -2H+E_1+E_2+E_5
  +w_0+w_1+w_2+q_0+3q_1,\\
b_9={}&
  -H+w_0+w_1+w_2+q_0+3q_1,\\
b_{10}={}&
  -3H+E_1+E_2+E_3+2E_4+E_5
  +w_0+w_1+w_2+3q_0+q_1,\\
b_{11}={}&
  3H-E_1-E_2-E_3-E_4-E_5
  -w_0-w_1-w_2-2q_0-2q_1,
\end{align*}
the permutation action of $\mathfrak g$ is visible by expressing
$$
g_1=(b_1,b_2)(b_3,b_5)(b_4,b_6)
  (b_7,b_8)(b_9,b_{10}),
  $$
  $$
  g_2=(b_1,b_6,b_3)(b_2,b_5,b_4)(b_8,b_9).
$$
\end{proof}
%This yields effective criteria for stable rationality: 
\begin{coro}\label{coro:dp4stab}
    Let $S$ be a quartic del Pezzo surface over a field $k$ of characteristic zero such that $S(k)\ne \varnothing$. Then $S$ is stably rational over $k$ if and only if one of the following equivalent conditions holds:
    \begin{enumerate}
           \item $\Pic(\bar{S})$ is a stably permutation $\mathfrak g$-module,
           \item $S$ satisfies Condition {\bf(H1)} of (1.1),
\item $S$ is not $k$-minimal or the $\mathfrak g$-action on $\Pic(\bar{S})$ factors through one of the four subgroups in Proposition~\ref{prop:h1}.
     \end{enumerate}
\end{coro}
\begin{proof}
    The equivalence of the conditions follows from Proposition~\ref{prop:h1} and Lemma~\ref{lemm:stablyperm}. Corollary~\ref{coro:rattorsor} and Proposition~\ref{prop:bctsmain} show that (1) implies the stable rationality of $S$. The converse is in \cite{Maninperfect}.
\end{proof}

\begin{prop}
    \label{prop:cubictypeI3}
    The smooth cubic hypersurface $X\subset\bP^{r+4}$ given by
    $$
    x_4(x_1^2+2x_1x_2)+x_3(x_1^2+x_1x_2+x_2^2)+x_3^3-x_3x_4^2+x_4^3+2x_5^3+\sum_{i=1}^rw_i^3=0
    $$
    is stably rational over $k=\bQ$, for any $r\geq 0.$
\end{prop}

\begin{proof}
  One can directly check that
$X$ is smooth when $r=0$. Adding the Fermat cubic in variables $w_i$ preserves the smoothness for every $r\geq 1$. Consider the projection 
    $$
    X\dashrightarrow \bP^{r+1},\quad (x_i,w_j)\to (x_4,x_5,w_1,\ldots,w_r).
    $$
    Let $K=k(\bP^{r+1})=k(a,b_1,\ldots,b_r)$. Set $a=x_4/x_5$, and $b_i=w_i/x_5$. Renaming the coordinates 
    $$
    Y_1=x_1,\quad Y_2=x_2, \quad Y_3=x_5,\quad Y_4=x_3,
    $$
 the generic fiber $X_\eta\subset\bP^3_{Y_1,Y_2,Y_3,Y_4}$ is given by
$$
Y_3(aY_1^2+2aY_1Y_2+(a^3+\beta)Y_3^2)+Y_4(Y_1^2+Y_1Y_2+Y_2^2-a^2Y_3^2+Y_4^2)=0,
$$
where $\beta=2+\sum_{i=1}^rb_i^3$. One can check that $X_\eta$ is smooth over $K$. Contracting the line $l_\eta=\{Y_3=Y_4=0\}\subset\bP^3$ yields a quartic del Pezzo surface $S$ birational to $X_\eta$ over $K$. Consider the  conic bundle obtained via projection from $l_\eta$
$$
\pi_\eta:X_\eta\to\bP^1,\quad (Y_1,Y_2,Y_3,Y_4)\to (Y_3,Y_4).
$$
Over the affine chart $t=Y_4/Y_3$, the generic fiber is the conic given by 
$$
f(Y_1,Y_2):=(t + a)Y_1^2 + (t + 2a)Y_1Y_2 + tY_2^2 + t^3 -a^2t + a^3 + \beta=0.
$$
The discriminant of the quadratic form $f$ is 
$$
\mathrm{disc}(f)=-\frac14(t^3-a^2t+a^3+\beta)(3t^2-4a^2).
$$
Let $c_1(t)=t^3-a^2t+a^3+\beta$, and $\rho_1,\rho_2,\rho_3$ its three roots. One sees that
$$
\Delta:=\disc(c_1)=-23a^6 - 54a^3\beta - 27\beta^2
$$
is not a square in $K$. Thus $\Gal(K_1/K)=\fS_3$, where $K_1=K(\rho_1,\rho_2,\rho_3)$. Put 
$$
d_1=\rho_2-\rho_3,\quad d_2=\rho_3-\rho_1,\quad d_3=\rho_1-\rho_2.
$$
The coefficients of $c_1(t)$ imply that 
$$
d_i^2=4a^2-3\rho_i^2.
$$
The singular fibers over $\rho_i$ split over $K_1$, and their components are
\begin{align}\label{eqn:I3fiber1}
F_i^\pm:    2(\rho_i+a)Y_1+(\rho_i+2a\pm d_i)Y_2=0,\quad i=1,2,3.
\end{align}
Let $c_2(t)=3t^2-4a^2$. The splitting field of $c_2$ is $K(\sqrt{3})$, and $c_2$ has two roots 
$\rho_4=\frac{2a}{\sqrt3}$ and $\rho_5=-\frac{2a}{\sqrt3}$. We check that $\sqrt{3}\notin K_1$. Indeed, $K(\sqrt 3)$ is different from the unique quadratic subextension $K(\sqrt{\Delta})/K$ of $K_1/K$, since ${\Delta/3}$ is  not a square in $K$. Set $K_2=K_1(\sqrt 3)$. It follows that
$$
\Gal(K_2/K)=\Gal(K(\sqrt 3)/K)\times \Gal(K_1/K)=\fC_2\times\fS_3.
$$
Over $K_2$, each singular fiber of $\pi_\eta$ is defined. Choose $D_1\in\bar K$ such that 
$$
D_1^2=(-32\sqrt3-52)a^4-(24\sqrt3+36)a\beta,
$$
and put 
\begin{align}
    \label{eqn:I3D2}
D_2=\frac{4a}{D_1}d_1d_2d_3.
\end{align}
One can check that 
\begin{align*}
D_2^2=(32\sqrt3-52)a^4+(24\sqrt3-36)a\beta.
\end{align*}
The singular fibers over $\rho_4$ and $\rho_5$ split over $K_2(D_1)$, with components: 
\begin{align}\label{eqn:I3fiber2}
    F_4^\pm&: 2a(3+2\sqrt{3} )Y_1+a(6+2\sqrt3)Y_2\pm{D_1}=0,\\\notag
    F_5^\pm&: 2a(3-2\sqrt{3} )Y_1+a(6-2\sqrt3)Y_2\pm{D_2}=0.
\end{align}
Thus, the splitting field of $X_\eta$ and $S$ is 
$$
K_2(D_1)=K(\sqrt3,\rho_1,\rho_2,\rho_3)(D_1)=K_1(D_1).
$$
We show that $D_1\notin K_2$. 
%Note that $D_1\notin K(\sqrt3)$, since the norm
%$$
%N_{K(\sqrt3)/K}(D_1^2)=D_1^2D_2^2=16a^2\Delta
%$$
%is not a square in $K$. 
One can check that the minimal polynomial of $D_1$ over $K$ is 
\begin{align}
 \label{eqn:I3D1mini}   
 (t^2-D_1^2)(t^2-D_2^2)\in K[t].
\end{align}
Then $K(D_1)/K$ has degree 4 and is not Galois since $D_2\not\in K(D_1)$.
However, any index 4 subgroup in $\Gal(K_2/K)=\fC_2\times\fS_3$ is normal, which implies that $D_1\notin K_2$. It follows that 
$$
[K_1(D_1):K]=[K_2(D_1):K]=[K_2(D_1):K_2][K_2:K]=24.
$$ 

%It follows that  $D_1=\lambda\sqrt\Delta$, for some $\lambda\in K(\sqrt3)$, which is impossible since the norm 
%$$
%N_{K(\sqrt3)/K}({D_1^2}/{\Delta})=\frac{16a^2}{\Delta}
%$$
%is not a square in $K$. 
 Let $G=\Gal(K_1(D_1)/K)$. From~\eqref{eqn:I3D2}, one sees that $K_2(D_1)$ is the splitting field of \eqref{eqn:I3D1mini} over $K_1$. It follows that $K_1(D_1)/K_1$ is Galois, and $\Gal(K_1(D_1)/K_1)=\fC_2^2$, since it has 3 different quadratic subextensions $K_1(\sqrt3)$, $K_1(D_1+D_2)$ and $K_1(D_1-D_2)$. We determine the image of $\Gal(K_1(D_1)/K_1)$ in $W(\sD_5)$ by identifying its action on~\eqref{eqn:I3fiber2}.
\begin{itemize}
    \item One generator of $\Gal(K_1(D_1)/K_1)$ swaps $D_1$ and $D_2$, and maps to $(4,5)\in W(\sD_5)$.
    \item The other generator of $\Gal(K_1(D_1)/K_1)$ changes the signs of $D_1$ and $D_2$ simultaneously, and maps to $c_4c_5\in W(\sD_5).$
\end{itemize}
The subextension 
$$
K_1(D_1)\supset K_1\supset K
$$
yields an extension of Galois groups
\begin{align}
    \label{eqn:I3ext}
    1\to \fC_2^2\to G\to\fS_3\to 1.
\end{align}
The extension~\eqref{eqn:I3ext} splits since we know that 
$
  K_1\cap K(D_1)=K.  
$
Any element in $\Gal(K_1/K)=\fS_3$ can be lifted to $G$ as an element fixing $D_1$. We describe their actions on singular fibers~\eqref{eqn:I3fiber1} and~\eqref{eqn:I3fiber2}.
\begin{itemize}
    \item A lift of $(1,2,3)\in\fS_3$ to $G$ permutes $\rho_1\mapsto\rho_2\mapsto\rho_3$, and fixes $D_1$ and $D_2$. Its image in $W(\sD_5)$ is $(1,2,3)$.
    \item A lift $\tau$ of $(2,3)\in\fS_3$ fixes $\rho_1$ and $D_1$, and swaps $\rho_2$ and $\rho_3$. One can check that $\tau (d_1)=-d_1$, $\tau(d_2)=-d_3$, $\tau(d_3)=-d_2$, and thus $\tau(D_2)=-D_2$. Its image in $W(\sD_5)$ is $c_1c_2c_3c_5(2,3)$.
\end{itemize}
Thus the image of $\mathfrak g_K$ in $W(\sD_5)$ is 
$$
\langle (4,5),c_4c_5,(1,2,3),c_1c_2c_3c_5(2,3)\rangle,
$$
which is conjugate to the group of type $I_3$ listed in Proposition~\ref{prop:h1}. Applying Corollary~\ref{coro:dp4stab}, we conclude that $X$ is stably rational over $\bQ$. 
\end{proof} 

\begin{thebibliography}{99}
\bibitem[CTSans]{CTSansucDuke} J.-L. Colliot-Th\'{e}l\`ene and J.-J. Sansuc. \newblock La descente sur les vari\'{e}t\'{e}s rationnelles. {II}. \newblock {\em Duke Math. J.}, 54(2):375--492, 1987.
\bibitem[KSS]{KSS} B.~Kunyavski\u~i, A.~Skorobogatov, and M.~Tsfasman. \newblock del {P}ezzo surfaces of degree four. \newblock {\em M\'em. Soc. Math. France (N.S.)}, (37):113, 1989.
\bibitem[Maninp]{Maninperfect} Yu. \~I. Manin. \newblock Rational surfaces over perfect fields. \newblock {\em Inst. Hautes \'Etudes Sci. Publ. Math.}, (30):55--113, 1966.
\bibitem[TY]{TY} Yu. Tschinkel and K.~Yang. \newblock Potentially stably rational del {P}ezzo surfaces over nonclosed fields. \newblock In {\em Combinatorial and additive number theory. {III}}, volume 297 of {\em Springer Proc. Math. Stat.}, pages 227--233. Springer, Cham, 2020.
\bibitem[bctsps]{bctspstably} A.~Beauville, J.-L. Colliot-Th\'el\`ene, J.-J. Sansuc, and P.~Swinnerton-Dyer. \newblock Vari\'et\'es stablement rationnelles non rationnelles. \newblock {\em Ann. of Math. (2)}, 121(2):283--318, 1985.
\bibitem[manin-]{manin-book} Yu. \~I. Manin. \newblock {\em Cubic forms}, volume~4 of {\em North-Holland Mathematical Library}. \newblock North-Holland Publishing Co., Amsterdam, second edition, 1986.
\end{thebibliography}
\end{document}
